\section{Conclusion}
\label{sec:conclusion}

This work conducted a controlled behavioral study of trajectory supervision for temporal state tracking in vision-language models. Using Qwen2.5-VL-3B-Instruct on a card-relocation task with discrete, observable state changes, we compared trajectory-supervised training (Model B) against a format-matched corrupted-trajectory control (B$'$v3).

\textbf{Primary finding}: On a held-out 10-video test set, trajectory-supervised Model B achieved 20\% final-position accuracy, statistically indistinguishable from the corrupted control B$'$v3 (also 20\%, McNemar's $p=1.00$). Per-round accuracy was numerically higher for B (25\% vs. 18\%) but not significant ($p=0.34$). Neither model achieved a single complete six-round exact match. Native video diagnostics confirmed true temporal input processing (24 frames per video), yet the model achieved 0\% complete trace accuracy and \textasciitilde19\% single-round accuracy.

\textbf{Evidence for format learning}: 
\begin{enumerate}
    \item Model B performs worse (10\%) on final-answer prompts despite learning trajectories, suggesting format-specific learning
    \item Corrupted control B$'$v3 achieves comparable accuracy despite invalid labels, indicating output distribution learning
    \item B$'$v3 outperforms B on move-accuracy, directly contradicting visual-grounding hypothesis
    \item Zero exact matches for either model suggest local prediction rather than persistent state maintenance
\end{enumerate}

\textbf{Methodological contributions}:
\begin{enumerate}
    \item \textbf{Format-matched negative control}: B$'$v3 construction via cross-example Hungarian move-count matching demonstrates how to create valid negative controls that preserve statistical properties while breaking visual correspondence
    \item \textbf{Dataset validity auditing}: Discovery of target $\equiv$ initial\_position confound across all 50 videos demonstrates the necessity of ground-truth auditing
    \item \textbf{Small-sample statistical rigor}: Exact binomial CIs, McNemar's test, power analysis, and transparent uncertainty quantification provide a template for small-$n$ VLM evaluation
    \item \textbf{Complete reproducibility}: 24 journey files documenting experimental iterations, seed 20260917 fixed split, complete artifact verification table, and hardcoded hyperparameters enable full reproduction
\end{enumerate}

\textbf{Limitations}:
\begin{itemize}
    \item Sample size: $n=10$ with MDE $\approx 60\%$ limits precision; replication with $n \geq 30$ needed
    \item Single model: Qwen2.5-VL-3B-Instruct only; findings may not generalize to larger models or different architectures
    \item Synthetic task: Card-relocation is artificially simple; naturalistic videos may show different patterns
    \item Frozen vision tower: Only language-layer fine-tuning; full-model tuning might differ
\end{itemize}

\textbf{Interpretation}:
We frame this as a \emph{cautious null result and controlled pilot study}, not a claim that trajectory supervision \emph{never} works or that VLMs are \emph{incapable} of tracking. Rather, we report that under these specific conditions---Qwen2.5-VL-3B, frozen vision tower, 40 training examples, this task---trajectory supervision did not demonstrate clear, consistent advantage over format-matched corrupted supervision. This conditional finding invites replication with:
\begin{enumerate}
    \item Larger test sets ($n \geq 30$)
    \item Larger models (7B, 72B variants)
    \item Diverse temporal reasoning tasks
    \item Unfrozen vision encoders
    \item Mechanistic interpretability analysis (attention, probing)
\end{enumerate}

\textbf{Broader implications}:
This work demonstrates the necessity of rigorous controls and validity auditing in VLM evaluation. Format-matched negative controls are feasible and informative. Small-sample null results, when properly reported with statistical uncertainty and audit trails, constitute legitimate scientific findings that can guide future research and prevent wasteful exploration of ineffective approaches.

\textbf{Future directions}:
\begin{enumerate}
    \item Scale to larger models and test sets to achieve adequate statistical power
    \item Compare trajectory supervision to other intermediate-supervision approaches (action descriptions, object annotations)
    \item Test on naturalistic videos to assess task-specificity of findings
    \item Employ mechanistic probes (attention visualization, representation geometry) to understand what models learn
    \item Investigate architectural modifications (explicit memory, recurrent video encoders) that might enable state tracking
\end{enumerate}

In conclusion, this controlled study provides evidence that trajectory supervision, under these conditions, teaches output format more readily than visual-temporal semantics. We hope this work serves as both a methodological template and a cautionary note for future VLM evaluation studies: rigorous controls and transparent null reporting are more valuable than inflated effect claims.
